% GENERATED FILE. Edit canonical lessons, then run npm run generate:books.
% canonical-source-hash: fnv1a64:6829eb68d3171468
% canonical-lessons: ZH-C374-renkou, ZH-C374-zhengfu, ZH-C374-wenhua, ZH-C374-minzu, ZH-C374-fazhan

\chapter{People and Government}
\label{ch:zh-people-and-government}
\hlchaptermodality{374}

\begin{chapteropening}
\textbf{By the end of this chapter:} \emph{I can talk about a population, the government, culture, ethnic groups and how a place develops.}

Complete the last lesson of chapter 374: I can talk about a population, the government, culture, ethnic groups and how a place develops.
\end{chapteropening}

\section[\texorpdfstring{\zh{人口} (rénkǒu)}{人口 (rénkǒu)}]{\zh{人口} (rénkǒu) — population}

\label{lesson:ZH-C374-renkou}



\begin{quote}
Before the new one: say the Mandarin for naive, then the Mandarin for to take it seriously.
\end{quote}

\subsection*{You'll want to know: \zh{人口}}
\textbf{\zh{人口}} — \emph{rénkǒu} — \textquotedblleft{}population\textquotedblright{}.

\begin{quote}
\emph{zhège chéngshì de rénkǒu yǒu yìqiān wàn} — This city has a population of ten million.
\end{quote}

\begin{cousinweb}
\textbf{\zh{人口}} (\emph{rénkǒu}) is built from \zh{人} 'a person' and \zh{口} 'a mouth', both inherited from Middle Chinese, \emph{nyin} and \emph{khuwX} (Guangyun, in Baxter's transcription), which glosses the second 'what a person speaks and eats with'. Counting the mouths of a place to mean its people is not given there.
\end{cousinweb}

\begin{grammarlens}[title={What you've built}]
One more word for giving reasons.
\end{grammarlens}

\subsection*{Your turn}
\begin{itemize}
\raggedright
  \item \emph{Say it:} \emph{rénkǒu}
  \item \emph{Say it:} \emph{rénkǒu}, once more
  \item \emph{Say it:} \emph{zhège chéngshì de rénkǒu yǒu yìqiān wàn} — This city has a population of ten million.
  \item \emph{Recall:} say the Mandarin for naive, then the Mandarin for to take it seriously, then say \emph{rénkǒu} again
\end{itemize}

\begin{tcolorbox}[breakable,colback=teal!4,colframe=teal!35!black,title={Before you move on}]
What does \zh{人口} mean? (\textquotedblleft{}Population\textquotedblright{}.) Say it once more.
\end{tcolorbox}
\section[zhèngfǔ]{zhèngfǔ — government}

\label{lesson:ZH-C374-zhengfu}



\begin{quote}
Before the new one: say the Mandarin for to take it seriously, then the Mandarin for the population (of a place).
\end{quote}

\subsection*{You'll want to know: zhèngfǔ}
\emph{zhèngfǔ} — \textquotedblleft{}government\textquotedblright{}.

\begin{quote}
\emph{zhèngfǔ yīnggāi guānxīn lǎo rén} — The government should look after old people.
\end{quote}

Its characters come later; for now it is a word for the ear.

\begin{cousinweb}
\emph{zhèngfǔ} is built from \emph{zhèng} 'governing' and \emph{fǔ} 'an office', both inherited from Middle Chinese, \emph{tsyengH} and \emph{pjuX} (Guangyun, in Baxter's transcription), which glosses the first 'to set right' and the second 'an office of state' and 'a store of documents'.
\end{cousinweb}

\begin{grammarlens}[title={What you've built}]
One more word for giving reasons.
\end{grammarlens}

\subsection*{Your turn}
\begin{itemize}
\raggedright
  \item \emph{Say it:} \emph{zhèngfǔ}
  \item \emph{Say it:} \emph{zhèngfǔ}, once more
  \item \emph{Say it:} \emph{zhèngfǔ yīnggāi guānxīn lǎo rén} — The government should look after old people.
  \item \emph{Recall:} say the Mandarin for to take it seriously, then the Mandarin for the population (of a place), then say \emph{zhèngfǔ} again
\end{itemize}

\begin{tcolorbox}[breakable,colback=teal!4,colframe=teal!35!black,title={Before you move on}]
What does zhèngfǔ mean? (\textquotedblleft{}Government\textquotedblright{}.) Say it once more.
\end{tcolorbox}
\section[wénhuà]{wénhuà — culture}

\label{lesson:ZH-C374-wenhua}



\begin{quote}
Before the new one: say the Mandarin for the population (of a place), then the Mandarin for the government.
\end{quote}

\subsection*{You'll want to know: wénhuà}
\emph{wénhuà} — \textquotedblleft{}culture\textquotedblright{}.

\begin{quote}
\emph{wǒ hěn xǐhuan Zhōngguó wénhuà bǐrú hànzì} — I really like Chinese culture, the characters for one.
\end{quote}

Its characters come later; for now it is a word for the ear.

\begin{cousinweb}
\emph{wénhuà} is built from \zh{文} 'writing' and \emph{huà} 'to transform', both inherited from Middle Chinese, \emph{mjun} and \emph{xwaeH} (Guangyun, in Baxter's transcription), which glosses the second 'to change' and 'to transform by example'. The sense 'culture' is not given there.
\end{cousinweb}

\begin{grammarlens}[title={What you've built}]
One more word for giving reasons.
\end{grammarlens}

\subsection*{Your turn}
\begin{itemize}
\raggedright
  \item \emph{Say it:} \emph{wénhuà}
  \item \emph{Say it:} \emph{wénhuà}, once more
  \item \emph{Say it:} \emph{wǒ hěn xǐhuan Zhōngguó wénhuà bǐrú hànzì} — I really like Chinese culture, the characters for one.
  \item \emph{Recall:} say the Mandarin for the population (of a place), then the Mandarin for the government, then say \emph{wénhuà} again
\end{itemize}

\begin{tcolorbox}[breakable,colback=teal!4,colframe=teal!35!black,title={Before you move on}]
What does wénhuà mean? (\textquotedblleft{}Culture\textquotedblright{}.) Say it once more.
\end{tcolorbox}
\section[mínzú]{mínzú — ethnic group; nationality}

\label{lesson:ZH-C374-minzu}



\begin{quote}
Before the new one: say the Mandarin for the government, then the Mandarin for a people's culture.
\end{quote}

\subsection*{You'll want to know: mínzú}
\emph{mínzú} — \textquotedblleft{}ethnic group; nationality\textquotedblright{}.

\begin{quote}
\emph{Zhōngguó yǒu wǔshíliù ge mínzú} — China has fifty-six ethnic groups.
\end{quote}

Its characters come later; for now it is a word for the ear.

\begin{cousinweb}
\emph{mínzú} is built from \emph{mín} 'the people' and \emph{zú} 'a clan', both inherited from Middle Chinese, \emph{mjin} and \emph{dzuwk} (Guangyun, in Baxter's transcription), which glosses the first 'the multitude' and the second 'a lineage'.
\end{cousinweb}

\begin{grammarlens}[title={What you've built}]
One more word for giving reasons.
\end{grammarlens}

\subsection*{Your turn}
\begin{itemize}
\raggedright
  \item \emph{Say it:} \emph{mínzú}
  \item \emph{Say it:} \emph{mínzú}, once more
  \item \emph{Say it:} \emph{Zhōngguó yǒu wǔshíliù ge mínzú} — China has fifty-six ethnic groups.
  \item \emph{Recall:} say the Mandarin for the government, then the Mandarin for a people's culture, then say \emph{mínzú} again
\end{itemize}

\begin{tcolorbox}[breakable,colback=teal!4,colframe=teal!35!black,title={Before you move on}]
What does mínzú mean? (\textquotedblleft{}Ethnic group; nationality\textquotedblright{}.) Say it once more.
\end{tcolorbox}
\section[fāzhǎn]{fāzhǎn — to develop; to expand}

\label{lesson:ZH-C374-fazhan}



\begin{quote}
Before the new one: say the Mandarin for a people's culture, then the Mandarin for an ethnic group.
\end{quote}

\subsection*{You'll want to know: fāzhǎn}
\emph{fāzhǎn} — \textquotedblleft{}to develop; to expand\textquotedblright{}.

\begin{quote}
\emph{zhège chéngshì fāzhǎn de hěn kuài suǒyǐ fángzi hěn guì} — This city is growing fast, so houses are expensive.
\end{quote}

Its characters come later; for now it is a word for the ear.

\begin{cousinweb}
\emph{fāzhǎn} is built from \emph{fā} 'to set going' and \emph{zhǎn} 'to open out', both inherited from Middle Chinese, \emph{pjot} and \emph{trjenX} (Guangyun, in Baxter's transcription), which glosses each with the same word, 'to unfold'.
\end{cousinweb}

\begin{grammarlens}[title={What you've built}]
One more word for giving reasons.
\end{grammarlens}

\subsection*{Your turn}
\begin{itemize}
\raggedright
  \item \emph{Say it:} \emph{fāzhǎn}
  \item \emph{Say it:} \emph{fāzhǎn}, once more
  \item \emph{Say it:} \emph{zhège chéngshì fāzhǎn de hěn kuài suǒyǐ fángzi hěn guì} — This city is growing fast, so houses are expensive.
  \item \emph{Recall:} say the Mandarin for a people's culture, then the Mandarin for an ethnic group, then say \emph{fāzhǎn} again
\end{itemize}

\begin{tcolorbox}[breakable,colback=teal!4,colframe=teal!35!black,title={Before you move on}]
What does fāzhǎn mean? (\textquotedblleft{}To develop; to expand\textquotedblright{}.) Say it once more.
\end{tcolorbox}
